\documentclass[aspectratio=169]{beamer}
\usetheme{metropolis}
\usepackage{nexus_theme}

\title{Anatomy of a Research Paper}
\subtitle{Module 0.1: How to Read Like a Scientist}
\author{Nexus Scholar Suite --- Modern Research Methodologies}
\date{}

\begin{document}

% --- TITLE SLIDE ---
\begin{frame}[plain]
  \titlepage
\end{frame}

% --- LEARNING OBJECTIVES ---
\begin{frame}{What You'll Learn}
  By the end of this lesson, you will be able to:
  \begin{enumerate}
    \item \textbf{Identify} the standard structural components of an academic paper.
    \item \textbf{Explain} the specific purpose of the IMRAD sections.
    \item \textbf{Apply} strategic reading techniques (the 3-Pass Approach).
    \item \textbf{Distinguish} between primary articles, reviews, and preprints.
  \end{enumerate}
\end{frame}

% --- THE PROBLEM ---
\begin{frame}{The Novice's Mistake}
  \begin{columns}[T]
    \begin{column}{0.48\textwidth}
      \begin{alertblock}{Reading Like a Novel}
        \begin{itemize}
          \item Starting at word 1, ending at the conclusion.
          \item Getting stuck in the Methods section.
          \item Spending 3 hours on a paper that isn't relevant.
          \item Exhaustion and burnout.
        \end{itemize}
      \end{alertblock}
    \end{column}
    \begin{column}{0.48\textwidth}
      \begin{block}{Reading Like a Scientist}
        \begin{itemize}
          \item Reading strategically based on goals.
          \item Treating the paper like an encyclopedia.
          \item Triaging in 10 minutes.
          \item Knowing exactly where to look.
        \end{itemize}
      \end{block}
    \end{column}
  \end{columns}
\end{frame}

% --- IMRAD ---
\begin{frame}{The Blueprint: IMRAD}
  \begin{exampleblock}{Introduction, Methods, Results, And Discussion}
    \begin{itemize}
      \item \textbf{Introduction:} The Porch (Why are we doing this?)
      \item \textbf{Methods:} The Foundation (How did we do it?)
      \item \textbf{Results:} The Rooms (What did we find?)
      \item \textbf{Discussion:} The Neighborhood (What does it mean?)
    \end{itemize}
  \end{exampleblock}
  \vspace{0.3cm}
  \textit{Note: The Abstract is a miniature version of the entire IMRAD structure.}
\end{frame}

% --- THE 3-PASS APPROACH ---
\begin{frame}{The Three-Pass Approach (Keshav, 2007)}
  \begin{center}
  \resizebox{0.92\linewidth}{!}{
  \begin{tikzpicture}[node distance=1.2cm and 1.8cm]
    \node[card, fill=NexusBlue!10, draw=NexusBlue, text width=3.2cm] (pass1) {
      \textbf{Pass 1: Triage}\\
      5-10 minutes\\
      Title, Abstract, Intro, Conclusion\\
      \textit{Goal: Is it relevant?}};
      
    \node[primarycard, right=of pass1, text width=3.2cm] (pass2) {
      \textbf{Pass 2: The Core}\\
      1 hour\\
      Figures, Tables, Results, Discussion\\
      \textit{Goal: Grasp findings}};
      
    \node[successcard, right=of pass2, text width=3.2cm] (pass3) {
      \textbf{Pass 3: Deep Dive}\\
      1-5 hours\\
      Methods (Line by line)\\
      \textit{Goal: Re-implement / Critique}};

    \draw[flowarrow] (pass1) -- node[above, font=\scriptsize\bfseries] {If Yes} (pass2);
    \draw[accentarrow] (pass2) -- node[above, font=\scriptsize\bfseries] {If Critical} (pass3);
  \end{tikzpicture}
  }
  \end{center}
\end{frame}

% --- TYPES OF PAPERS ---
\begin{frame}{Types of Academic Literature}
  \begin{block}{Not all papers are created equal}
    \begin{itemize}
      \item \textbf{Primary Research:} Original data and experiments (Uses IMRAD).
      \item \textbf{Review Articles:} Synthesis of dozens/hundreds of primary papers (Often skips Methods, focuses on themes).
      \item \textbf{Preprints (e.g., arXiv):} Early drafts that have \textbf{not} been peer-reviewed. Read with extra skepticism.
    \end{itemize}
  \end{block}
\end{frame}

% --- KEY TAKEAWAY ---
\begin{frame}{Key Takeaway}
  \begin{block}{What We Accomplished}
    \begin{enumerate}
      \item \textbf{IMRAD}: You know the universal blueprint of science.
      \item \textbf{Non-linear reading}: You have permission to skip the Methods section.
      \item \textbf{Efficiency}: You can triage a paper in under 10 minutes.
    \end{enumerate}
  \end{block}
  \vspace{0.3cm}
  \textbf{Next:} Module 0.2 --- Research Methodologies
\end{frame}

% --- EXERCISE PREVIEW ---
\begin{frame}{Your Turn: Exercise}
  \begin{exampleblock}{Exercise: The 10-Minute Triage}
    Set a timer for 10 minutes. Open a dense research paper provided in the exercise files, and use Pass 1 to answer three specific questions without reading the main body text.
  \end{exampleblock}
\end{frame}

% --- STANDOUT CLOSE ---
\begin{frame}[standout]
  You don't need to read every word to understand the science.\\
  You just need to know where to look.
\end{frame}

\end{document}
